\section{Capítulo 7 --- Integração da tabela imperfeita à recomendação}

Uma secretaria analisa evasão em cursos técnicos. A base reúne matrícula, frequência e situação final, mas há duplicatas, ausências, mudanças de definição e grupos com exposição desigual.

\subsection{Questões objetivas}

\ItemHeader{1}{Objetiva}{grão}{Situação profissional e suporte quantitativo}

\TextoBase

A tabela deveria ter uma linha por estudante-semestre, mas transferências criam duas linhas para a mesma pessoa. Como a taxa de evasão usa matrículas elegíveis no denominador, a duplicação altera contagem e perfil do grupo. A secretaria precisa corrigir o grão antes de agregar. O parecer acadêmico informa que a transferência não equivale a uma segunda exposição ao risco: ela muda o vínculo administrativo dentro do mesmo semestre. A trilha de reconciliação deve conservar os dois registros de origem, selecionar o estado válido e provar unicidade na tabela analítica. O relatório será auditado por curso e turno, portanto uma exclusão aleatória poderia alterar justamente os estratos menores. A regra deverá usar chave composta, vigência e estado acadêmico, registrar conflitos e interromper a publicação quando mais de uma linha permanecer elegível para a mesma chave. Só depois desse teste o numerador de evasão e o denominador de estudantes expostos poderão ser agregados sob a mesma definição.

Essa definição institucional vigente também deve constar da linhagem publicada.

\FonteDidatica

\Comando Selecione a ação que restaura uma linha por estudante-semestre e protege numerador e denominador.

\begin{enumerate}[label=\Alph*)]

\item somar todas as linhas porque representam eventos reais

\item remover aleatoriamente uma duplicata

\item definir chave estudante-semestre, reconciliar transferências e testar unicidade antes da taxa

\item agregar por curso antes de investigar

\item usar a média das taxas publicadas

\end{enumerate}

\ItemHeader{2}{Objetiva}{taxa}{Situação profissional e suporte quantitativo}

\TextoBase

O relatório da secretaria registra no Quadro 1 evasões e estudantes elegíveis de dois cursos sob a mesma definição. A comparação deve usar taxas, sem classificar pela contagem absoluta nem ocultar o denominador menor.
\begin{DocumentBox}[Quadro 1 --- evasão por curso]
\begin{tabularx}{\textwidth}{@{}Yrr@{}}\toprule
\textbf{Curso} & \textbf{Evasões} & \textbf{Elegíveis}\\\midrule
A & 30 & 200\\ B & 18 & 80\\\bottomrule
\end{tabularx}
\end{DocumentBox}

\FonteDidatica

\Comando Calcule as taxas de A e B e selecione a comparação proporcional aos denominadores.

\begin{enumerate}[label=\Alph*)]

\item A=15\% e B=22,5\%; B tem maior taxa, com ressalva sobre denominador menor

\item A=30\% e B=18\%

\item A=12,5\% e B=7,5\%

\item A é pior porque tem mais casos

\item as taxas somam 37,5\%

\end{enumerate}

\ItemHeader{3}{Objetiva}{causalidade}{Situação profissional e suporte quantitativo}

\TextoBase

Cursos com mais atendimentos psicopedagógicos apresentam menor evasão. Entretanto, o atendimento foi priorizado para estudantes que solicitaram ajuda, e os cursos diferem em perfil e calendário. A secretaria avalia expansão, mas o desenho observacional ainda admite seleção e explicações alternativas.

\FonteDidatica

\Comando Selecione a conclusão que distingue associação observada de efeito causal e propõe evidência comparável.

\begin{enumerate}[label=\Alph*)]

\item concluir que atendimento causou toda a redução

\item concluir que atendimento aumenta risco

\item remover estudantes atendidos

\item tratar como associação sujeita a seleção; propor desenho comparável antes de atribuir efeito

\item usar correlação como prova causal

\end{enumerate}

\ItemHeader{4}{Objetiva}{ausência}{Situação profissional e suporte quantitativo}

\TextoBase

A frequência está ausente em 40 de 400 registros; 30 dessas ausências concentram-se em um campus cujo leitor falhou. Substituir vazios por zero classificaria falha de captura como falta. A secretaria precisa quantificar completude e investigar o mecanismo antes de imputar.

\FonteDidatica

\Comando Calcule a proporção ausente e selecione o tratamento compatível com a concentração por campus.

\begin{enumerate}[label=\Alph*)]

\item 1\% e ausência completamente aleatória

\item 10\% ausente e mecanismo dependente do campus; investigar processo antes de imputar

\item 40\% e excluir toda a coluna

\item 7,5\% e substituir por zero

\item 10\% e imputar média sem registrar

\end{enumerate}

\ItemHeader{5}{Objetiva}{ética}{Situação profissional e suporte quantitativo}

\TextoBase

Um ranking público detalha evasão por curso, turno e faixa etária. Algumas células possuem menos de cinco estudantes e, combinadas com informações locais, permitem inferir situações individuais. A transparência precisa ser equilibrada com finalidade, agregação e risco de reidentificação.

\FonteDidatica

\Comando Selecione a publicação que reduz reidentificação sem eliminar a finalidade analítica.

\begin{enumerate}[label=\Alph*)]

\item publicar porque nomes não aparecem

\item trocar nomes por códigos e manter células

\item adicionar cores para anonimizar

\item pedir ao LLM que decida a divulgação

\item agregar/suprimir células pequenas, limitar finalidade e documentar risco residual

\end{enumerate}

\ItemHeader{6}{Objetiva}{covariação}{Situação profissional e suporte quantitativo}

\TextoBase

Horas de trabalho e faltas apresentam covariância positiva e correlação $r=0,62$. Após verificar e retirar dois casos extremos em análise de sensibilidade, $r=0,28$. Os casos são legítimos; a secretaria precisa comunicar influência, tamanho e associação sem seleção oportunista nem alegação causal.

\FonteDidatica

\Comando Compare as duas correlações e selecione a interpretação que comunica sensibilidade sem excluir casos automaticamente.

\begin{enumerate}[label=\Alph*)]

\item a correlação aumentou 0,34

\item 62\% das faltas são causadas por trabalho

\item a associação linear cai 0,34 e é sensível aos extremos; relatar ambas e não inferir causa

\item covariância positiva prova unidade comum

\item retirar extremos sempre produz verdade

\end{enumerate}

\ItemHeader{7}{Objetiva}{auditoria}{Situação profissional e suporte quantitativo}

\TextoBase

O dashboard atual mostra evasão de 17\%, enquanto uma planilha anterior informa 14\% para período aparentemente igual. Não há registro de definição, janela, grão, filtros ou versão. Antes de escolher um valor, a secretaria precisa reconstruir as duas cadeias de cálculo.

\FonteDidatica

\Comando Selecione o procedimento que reconcilia definição, período, grão, filtros e versões antes da publicação.

\begin{enumerate}[label=\Alph*)]

\item publicar a média 15,5\%

\item reconciliar definição, janela, grão, filtros e versões antes de escolher um valor

\item usar o valor mais recente sem investigação

\item manter ambos sem rótulo

\item escolher o menor para evitar alarme

\end{enumerate}

\ItemHeader{8}{Objetiva}{risco relativo}{Situação profissional e suporte quantitativo}

\TextoBase

No mesmo período, o turno noturno registrou 48 evasões em 240 elegíveis; o diurno, 24 em 240. A secretaria quer comparar risco absoluto e razão entre taxas, mantendo explícito que o turno pode estar associado a outras condições não controladas.

\FonteDidatica

\Comando Calcule as duas taxas e sua razão e selecione a interpretação que preserva o limite causal.

\begin{enumerate}[label=\Alph*)]

\item diferença de 24 pontos percentuais

\item razão 0,5 com noturno protegido

\item noturno causa 100\% das evasões

\item taxas 20\% e 10\%; razão 2, associação que ainda requer controle de confundidores

\item taxas iguais porque denominadores coincidem

\end{enumerate}

\ItemHeader{9}{Objetiva}{recomendação}{Situação profissional e suporte quantitativo}

\TextoBase

Um campus adotou intervenção de permanência no mesmo mês em que o calendário acadêmico mudou; os demais não adotaram. A evasão caiu no campus participante. A secretaria precisa decidir o próximo passo sem atribuir toda a mudança à intervenção por simples comparação antes/depois.

\FonteDidatica

\Comando Selecione o desenho de continuidade que gera evidência comparável antes de uma expansão definitiva.

\begin{enumerate}[label=\Alph*)]

\item recomendar piloto ampliado com comparação e métricas prévias, não expansão causal definitiva

\item expandir imediatamente por correlação temporal

\item cancelar porque não há ensaio perfeito

\item ocultar mudança de calendário

\item comparar apenas antes/depois

\end{enumerate}

\ItemHeader{10}{Objetiva}{rastreabilidade}{Situação profissional e suporte quantitativo}

\TextoBase

A diretoria exige que cada taxa de evasão apresentada hoje possa ser refeita seis meses depois. Fontes recebem novas linhas, regras podem mudar e filtros do dashboard não ficam no PDF. O pacote de evidência precisa versionar toda a cadeia e seus responsáveis.

\FonteDidatica

\Comando Selecione o conjunto de artefatos necessário para reproduzir o indicador após mudanças de dados e regras.

\begin{enumerate}[label=\Alph*)]

\item guardar apenas PDF final

\item guardar prompt sem dados

\item fotografar o dashboard

\item registrar somente a ferramenta usada

\item versionar fonte, transformação, regra, filtro e artefato com responsável e data

\end{enumerate}

\subsection{Questões discursivas}

\ItemHeader{11}{Discursiva}{Integrar evidências e justificar decisão}{Caso profissional do capítulo}

\TextoBase

Uma base de evasão deveria possuir uma linha por estudante-semestre, mas transferências duplicam matrículas, frequência apresenta falhas por campus e regras mudaram durante o período. A recomendação final será usada para priorizar atendimento e precisa sobreviver a auditoria.

\FonteDidatica

\Comando Desenhe, em até 15 linhas, o fluxo da fonte à recomendação. Inclua teste de grão, ausência e domínio, transformação versionada, visual apropriado, reconciliação e evidência necessária para auditoria.

\ItemHeader{12}{Discursiva}{Integrar evidências e justificar decisão}{Caso profissional do capítulo}

\TextoBase

No mesmo período, o curso A registrou 30 evasões em 200 elegíveis e B, 18 em 80. O noturno registrou 48 em 240 e o diurno, 24 em 240. Os grupos diferem em tamanho e composição, e a secretaria considera concentrar recursos em apenas um deles.

\FonteDidatica

\Comando Calcule as quatro taxas, as diferenças em pontos percentuais e as razões A/B e noturno/diurno. Recomende uma prioridade provisória, distinguindo volume, risco relativo e limite causal.

\ItemHeader{13}{Discursiva}{Integrar evidências e justificar decisão}{Caso profissional do capítulo}

\TextoBase

Um painel público permite filtrar evasão por curso, turno e faixa etária. Em alguns recortes restam três ou quatro estudantes, e servidores locais conhecem a composição das turmas. A instituição precisa informar desempenho sem revelar situações individuais nem ampliar a finalidade original.

\FonteDidatica

\Comando Proponha uma política para células pequenas com limiar ou agregação justificada, finalidade, perfis de acesso, retenção e comunicação do risco residual. Inclua uma forma de testar reidentificação.

\ItemHeader{14}{Discursiva}{Integrar evidências e justificar decisão}{Caso profissional do capítulo}

\TextoBase

O dashboard informa evasão de 17\%, enquanto a planilha usada no mês anterior mostra 14\%. As equipes não conseguem localizar versão da definição, filtros, população elegível ou data de extração. A diretoria suspendeu a publicação até haver reconciliação.

\FonteDidatica

\Comando Crie um protocolo de reconciliação para 14\% e 17\%. Ordene verificações de definição, período, população, grão, filtros, transformação e versão; atribua responsáveis e declare a evidência para republicar.

\ItemHeader{15}{Discursiva}{Integrar evidências e justificar decisão}{Caso profissional do capítulo}

\TextoBase

Um campus adotou intervenção de permanência ao mesmo tempo em que alterou calendário e equipe. A evasão posterior caiu, mas não existe campus pareado nem métrica prévia de exposição à intervenção. A direção deseja expandir a política para toda a rede.

\FonteDidatica

\Comando Avalie a alegação causal. Identifique confundidores, proponha grupo ou período comparável, métricas prévias e uma regra mensurável de expansão que não dependa apenas da diferença antes/depois.

\newpage

\subsection{Respostas explicadas das questões pares}

\paragraph{Questão 2 --- resposta A.} Para A, $30/200=0{,}15=15\%$. Para B, $18/80=0{,}225=22{,}5\%$. Apesar de A ter mais casos absolutos, B apresenta taxa 7,5 pontos percentuais maior. A interpretação deve manter visível que B possui denominador menor; portanto, A é a resposta correta.
\[p_A=15\%,\qquad p_B=22{,}5\%\]

\paragraph{Questão 4 --- resposta B.} A ausência global é $40/400=0{,}10=10\%$. Desses 40 registros, 30 vêm do campus com falha de leitor, concentração incompatível com tratar a ausência como completamente aleatória. Preencher com zero confundiria falha de captura com falta; deve-se investigar e corrigir o processo antes de imputar. Logo, B.

\paragraph{Questão 6 --- resposta C.} A correlação passa de 0,62 para 0,28, queda de $0{,}62-0{,}28=0{,}34$. Isso mostra sensibilidade aos dois casos extremos, não licença para apagá-los. Como são legítimos, ambas as análises devem ser relatadas; nenhum dos coeficientes prova causalidade. Essa leitura está em C.

\paragraph{Questão 8 --- resposta D.} No noturno, $48/240=0{,}20=20\%$; no diurno, $24/240=0{,}10=10\%$. A razão é $20\%/10\%=2$: a taxa observada no noturno é o dobro. O cálculo descreve associação e ainda requer controle de perfil, trabalho e outras variáveis. Logo, D.

\paragraph{Questão 10 --- resposta E.} Guardar apenas o PDF não congela dados, regra ou recorte. A reprodução exige versão da fonte ou snapshot, transformação, definição da métrica, filtros e artefato, cada qual com responsável e data. Esse encadeamento permite explicar mudanças futuras e está completo somente na alternativa E.

\paragraph{Questão 12 --- critérios.} Esperam-se os cálculos $30/200=15\%$, $18/80=22{,}5\%$, $48/240=20\%$, $24/240=10\%$ e razão 2. A interpretação deve distinguir diferença absoluta, razão e tamanhos amostrais, além de propor comparação estratificada antes de atribuir efeito ao turno.

\paragraph{Questão 14 --- critérios.} A resposta deve reconciliar 17\% e 14\% por definição, janela, grão, filtros e versão; não vale publicar média entre indicadores incompatíveis. Deve propor um painel com contrato da métrica, denominadores, data de atualização, rastreabilidade e regra de bloqueio enquanto a divergência não for explicada.
